\chapter{Dry Eye Disease (Dry Eyes)}

\section*{Definition}
A multifactorial disease of the ocular surface characterized by loss of homeostasis of the tear film, accompanied by ocular discomfort symptoms and visual disturbance.

\section*{When to See a Doctor}
See your doctor if:
\begin{itemize}
  \item Persistent dry eye symptoms despite artificial tears, eye pain, photophobia, or vision changes
\end{itemize}

\section*{How It Develops}
Tear film instability from either reduced tear production (aqueous-deficient dry eye — Sjogren, age, medications) or excessive tear evaporation (meibomian gland dysfunction leading to lipid layer deficiency). Ocular surface inflammation creates a vicious cycle of tear dysfunction, hyperosmolarity, and epithelial damage.

\section*{Causes}
\begin{itemize}
  \item Aging (reduced tear production with age)
  \item Meibomian gland dysfunction (most common cause)
  \item Environmental factors (dry climate, AC, wind, screen time)
  \item Contact lens wear
  \item Sjogren syndrome (autoimmune)
  \item Medication side effects (antihistamines, decongestants, diuretics, antidepressants)
  \item Hormonal changes (menopause)
  \item Blepharitis (eyelid inflammation)
  \item Refractive surgery (LASIK)
  \item Systemic diseases (rheumatoid arthritis, lupus, diabetes)
\end{itemize}

\section*{Related Symptoms}
\begin{itemize}
  \item Eye redness
  \item Blurred vision
  \item Photophobia
  \item Eye fatigue
  \item Foreign body sensation
\end{itemize}

\section*{Medical Evaluation}
\begin{itemize}
  \item Your doctor will take a history of ocular symptoms, screen time, contact lens use, and medical conditions.
  \item A slit-lamp examination evaluates the tear film, cornea, conjunctiva, and eyelid margins.
  \item Tear break-up time (TBUT) and Schirmer test measure tear film stability and production.
  \item Meibomian gland assessment checks for gland dysfunction and obstruction.
  \item Blood tests for autoimmune markers (ANA, anti-Ro/SSA, anti-La/SSB) if Sjogren syndrome is suspected.
\end{itemize}

\section*{Treatment Options}
\begin{itemize}
  \item Artificial tears and lubricating gels are first-line for mild dry eye disease.
  \item Anti-inflammatory eye drops (cyclosporine, lifitegrast) for moderate to severe inflammatory dry eye.
  \item Meibomian gland expression, warm compresses, and lid hygiene for evaporative dry eye.
  \item Punctal plugs to reduce tear drainage in aqueous-deficient dry eye.
  \item Omega-3 fatty acid supplements may improve meibomian gland function and tear quality.
\end{itemize}

\section*{Self-Care and Home Management}
\begin{itemize}
  \item Use preservative-free artificial tears throughout the day, especially during screen use.
  \item Apply warm compresses to the eyelids for 5-10 minutes to unblock meibomian glands.
  \item Take regular breaks from screens using the 20-20-20 rule (every 20 minutes, look 20 feet away for 20 seconds).
  \item Use a humidifier in dry indoor environments to reduce tear evaporation.
  \item Avoid air conditioning or fan drafts directed at the eyes.
\end{itemize}

\section*{References}
\begin{thebibliography}{99}
\bibitem{dry_eye_disease:ref1} Craig JP, Nichols KK, Akpek EK, et al. ``TFOS DEWS II definition and classification report.'' Ocul Surf. 2017;15(3):276-283.
\bibitem{dry_eye_disease:ref2} Stapleton F, Alves M, Bunya VY, et al. ``TFOS DEWS II epidemiology report.'' Ocul Surf. 2017;15(3):334-365.
\bibitem{dry_eye_disease:ref3} Messmer EM. ``The pathophysiology, diagnosis, and treatment of dry eye disease.'' Dtsch Arztebl Int. 2015;112(5):71-81.
\bibitem{dry_eye_disease:ref4} Jones L, Downie LE, Korb D, et al. ``TFOS DEWS II management and therapy report.'' Ocul Surf. 2017;15(3):575-628.
\bibitem{dry_eye_disease:ref5} Lemp MA. ``Dry eye disease: a multi-factorial disease.'' In: Cornea. 4th ed. Elsevier; 2022.
\bibitem{dry_eye_disease:ref6} Foulks GN. ``Dry eye disease: clinical features and diagnosis.'' UpToDate. Wolters Kluwer; 2025.
\end{thebibliography}
